\documentclass[10pt,a4paper,twoside,bg=print,twocolumn,openany,nodeprecatedcode]{dndbook}
\usepackage[utf8]{inputenc}
\usepackage{tabulary}
\usepackage[normalem]{ulem}
\usepackage{color,soul}
\usepackage{float}
\usepackage{caption}
\usepackage{wrapfig}
\usepackage{tocloft}
\usepackage[toc]{multitoc}
\usepackage[hidelinks]{hyperref}
\raggedbottom
\extrafloats{2000}
\cftsetindents{section}{0em}{0em}
\cftsetindents{subsection}{0em}{0em}
\cftsetindents{subsubsection}{0.2em}{0em}
\setcounter{tocdepth}{1}
\renewcommand*{\multicolumntoc}{3}
\setlist[description]{nolistsep,listparindent=3pt,topsep=6pt}
\date{}
\begin{document}
\frontmatter
\tableofcontents
\mainmatter
\chapter{Feywild Overview}
\DndDropCapLine{D}{}omains of Delight are to the Feywild what Domains of Dread are to the Shadowfell: sequestered realms governed by powerful beings. Whereas a Domain of Dread is ruled by a Darklord, a Domain of Delight is ruled by an archfey—the most powerful of Fey creatures. An archfey gives form to their Domain of Delight, shaping it in ways unique to their personality. Some Domains of Delight are bright and cheery, while others are gloomy, but each one reflects the emotional state of its ruler. A Domain of Delight can be as small as a few acres or as big as a country.


This accessory helps you create Domains of Delight and the archfey who rule them, building on the information about the Feywild that appears in the Dungeon Master's Guide. The ideas, tips, and tables in this chapter are meant to spark your imagination. Use what excites and intrigues you, discard what doesn't, and make up the rest!

\section{Feywild Features}
The Feywild responds to unfettered emotion. It's not uncommon for flowers to turn and tremble if there's a heated argument nearby. If someone is filled with malice, their footprints might wither the grass under their feet or cause underground insects and worms to burrow to the surface. Birds chirp merrily in the presence of those who are joyous and squawk angrily at those who are perpetually dour. Nosy trees lean in to overhear whispers of conspiracy, eager for delicious tidbits they can gossip about later, and a rock might reshape itself to look like the creature that's happily sunning itself on the rock's surface.

Time and distance in the Feywild are mutable, as is the plane's geography. Roads are uncommon, and those that exist are as likely to change as the land around them. Because the distance between locations is not fixed and dilations in time are commonplace, a journey that took one hour yesterday might take three days tomorrow. Feywild natives are accustomed to the plane's mutability. For them, it's no more peculiar than the sun rising and setting on a Material Plane world.

Other features of the Feywild are described in the sections that follow. Think of them as a sampling of what the Feywild has to offer, for like the Material Plane, the Feywild is vast and diverse.

\subsection{Seelie and Unseelie Fey}
Seelie Fey and Unseelie Fey are two groups that often find themselves at odds. Seelie Fey cling to the trappings of civilization, value protocol, and uphold traditions. Unseelie Fey indulge their primal instincts, abhor adherence to protocol, and shun conformity. The two groups are not opposites morally or ethically; good and evil Fey can be found in both.

Seelie and Unseelie Fey gather in courts. The Seelie court is called the Summer Court, and the Unseelie court is called the Gloaming Court. Both courts stretch to the far corners of the Feywild, so their representatives can be encountered almost anywhere on this plane of existence. The Summer Court and the Gloaming Court are by no means the only great Fey courts, but they're the most wellknown to creatures on the Material Plane and the most widespread.

How the Summer Court and the Gloaming Court came to be is a mystery. Perhaps some Fey felt a deeper affinity with the natural world and chose to emulate it, while other Fey began to control nature, using magic to invent new ways of living. Whatever the case, innumerable Fey pursued these two paths, which became the two courts, and there have been squabbles between them ever since.

Each court tries to destabilize and demoralize the other. Both Fey courts have spies who dig up dirt, sow seeds of dissent, and cause mischief. Captured spies are either ransomed or made examples of in various humiliating, nonlethal ways.

Much of the gossip and chatter within the Feywild is fueled by the intrigue and drama between the two queens that rule the courts. Titania, the Summer Queen, is the regal and charismatic ruler of the Summer Court. Her court enjoys a tenuous peace with the unearthly ruler of the Gloaming Court, the Queen of Air and Darkness, who allows her kin to dabble in magic forbidden by the Seelie Court.

Conflicts between the Fey courts are often ritualized. Representatives of both courts gather in an amphitheater or field to have heated debates or energetic dance competitions that simulate combat, and these events are often laced with bawdy insults and lewd gestures. Only on rare occasions do things get physical, and even then, the Fey do little more than bite, scratch, and hurl mud at each other.

Every now and then, the two queens lock horns, causing tensions to rise throughout the Feywild. If only one of them can get her way, what would normally be a squabble can turn to bloodshed. But only the Fey queens can declare all-out war against each other, and it would take something extreme to have them abandon their relative peace and hurl their courts into violence and chaos.

The Summer Court and the Gloaming Court have nothing akin to a mutual defense pact, and the very thought of one is greeted by jeers on both sides. If a rising army of fomorians or some other threat endangers one court, it's widely understood that the other court will not intervene unless it, too, is threatened by the same enemy.

If you choose to align your domain and its archfey with one of the two courts, guidelines for how they're differentiated are given below.

\subsubsection{Summer Court (Seelie Fey)}

\begin{itemize}
\item Favors sunshine, butterflies, flowers, music, and singing birds
\item Values ceremony and refinement (for example, proper ways of speaking, formal etiquette, and extravagant dinner parties)
\item Favors elaborate, manufactured costuming with immaculate tailoring
\item Harshly judges those who don't exhibit the proper etiquette (adventurers who commit social blunders make fools of themselves, might be labeled buffoons, and might be laughed out of court)
\end{itemize}

\subsubsection{Gloaming Court (Unseelie Fey)}

\begin{itemize}
\item Favors gloom, twilight, cobwebs, fireflies, hooting owls, and croaking frogs
\item Values the intuitive and instinctual (for example, mystical rituals, visionaries, and firelit parties)
\item Shuns the constraints of civilization (instead wearing only unfinished natural materials and sleeping under the stars)
\item Dabbles in mysterious magic and rituals (adventurers can run afoul of curses if they don't follow the Unseelie ways)
\end{itemize}

\subsection{Spells in the Feywild}
Spells that manifest one way in the Material Plane might do so differently in the Feywild. A \textit{magic missile} spell might take the form of a giggling sprite that materializes next to the caster, fires off a barrage of tiny, glowing arrows, and disappears in a puff of sparkling fairy dust. Here are some other examples of how spells can be cosmetically reinterpreted in the Feywild:

\subparagraph{Crown of Madness}
The crown is made from gingerbread with icing filigree and candy gems.

\subparagraph{Find the Path}
A pixie-like spirit appears and guides the caster to the desired location. The spirit can't be harmed.

\subparagraph{Gust of Wind}
The wind carries the scent of flowers.

\subparagraph{Maze}
The demiplanar labyrinth created by the spell resembles a thorny hedge maze.

\subparagraph{Phantom Steed}
The steed looks like a giant, fuzzy caterpillar.

\subparagraph{Revivify}
A creature restored to life by this spell wakes up wondering if their entire life was all just a dream.

\subsection{Weather}
Strolling across a meadow or walking across a desert in the Feywild is often no different than doing the same on the Material Plane, but sometimes the Feywild needs to feel otherworldly. Playing with something as simple as the weather is an easy way to remind your players that their characters are not on the Material Plane.

The Feywild Weather table helps you determine fun weather effects. A weather effect such as this usually lasts no more than an hour or three.

\begin{DndTable}[header=Feywild Weather]{cX}

d8 & Weather Effect\\
1 & Flower blossoms rain from a sky filled with sparkling, pastel-colored clouds.\\
2 & A grand fairy drama plays in the sky as a stylistically rendered illusion or a swirling aurora.\\
3 & It rains tiny fruit tarts. They fall slowly and disappear just before they hit the ground. If caught from the air and eaten, they're delicious.\\
4 & Fog rolls in and plays harmless tricks on the characters, giving them fog mustaches, fog eyebrows, and fog wigs of many styles—perhaps even fog cloaks or fog companions.\\
5 & Snow begins to fall, and the flakes grow bigger over time. Flakes as big as dinner plates, wagon wheels, and even a waterwheel fall, but they're light as a feather.\\
6 & The sky fills with iridescent bubbles that lazily fall to the ground. Giggling pixie children chase and pop the bubbles but turn invisible and flee if spoken to or approached.\\
7 & A silvery rain lifts the spirits and brings a song to each heart that it touches.\\
8 & A howling wind blows through the party, and each party member must roll a d8. Anyone who rolls an 8 has a trinket or some other tiny, nonmagical item (such as a coin) stolen from them by weather spirits.\\
\end{DndTable}

\subsection{Domain Borders}
As an archfey's power waxes or wanes, their Domain of Delight can grow or shrink. If their domain grows big enough to abut or overlap another archfey's domain, a territorial dispute can arise. Until this dispute is resolved, other Fey denizens of the overlapping domains must defer to both rulers. Such disputes rarely last long; in the end, one archfey is given sufficient incentive to move elsewhere, or the two archfey learn to live with each other (and other denizens of the region must answer to both).

An archfey whose Domain of Delight doesn't overlap with another archfey's domain can surround their domain's border with walls of shimmering mist or some other magical effect that hides the domain from view and, if the archfey wishes, prevents creatures from entering or leaving the domain without a key, a guide, a password, the answer to a riddle, the archfey's consent, or magic.

\subsection{Feywild Guides}
The Feywild has its own illogical logic that visitors from other planes can never fully grasp. A guide who is well-versed in the ways of the Feywild can save travelers time and frustration, possibly even their freedom and their lives—by helping them avoid or get around the illogical logic. While visitors are easily vexed by roads that lead nowhere and forest trails that double back on themselves, a capable guide can see the road through the road and the trail beneath the trail. In doing so, the guide sidesteps the confusion and leads charges safely to their intended destination. Conversely, a bad guide can easily get an adventuring party into trouble.

Use the Feywild Guide Names, Feywild Guide Identities, and Feywild Guide Quirks tables to create Fey guides on the fly.

\begin{DndTable}[header=Feywild Guide Names]{cX}

d8 & Name\\
1 & Fetter\\
2 & Fiddlebones\\
3 & Moonray\\
4 & Pip\\
5 & Starlight\\
6 & Stumpwick\\
7 & Thistledown\\
8 & Whisperwind\\
\end{DndTable}

\begin{DndTable}[header=Feywild Guide Identities]{cX}

d8 & Identity\\
1 & Friendly forest gnome (use the \textbf{scout} stat block, but change its size to Small)\\
2 & Gloomy wood elf druid (use the \textbf{druid} stat block)\\
3 & Flamboyant \textbf{pixie}\\
4 & Happy-go-lucky \textbf{satyr}\\
5 & Ultra-competitive \textbf{sprite}\\
6 & Overly cautious \textbf{treant sapling} (see chapter 1 of The Wild Beyond the Witchlight for its stat block)\\
7 & Giggly goblin warlock (use the \textbf{cult fanatic} stat block, but change its size to Small)\\
8 & Humorless \textbf{centaur}\\
\end{DndTable}

\begin{DndTable}[header=Feywild Guide Quirks]{cX}

d8 & Quirk\\
1 & Always hearkens back and compares things to "the good old days"\\
2 & Can't tell a lie without wiggling their nose first\\
3 & Never takes off their hat\\
4 & Loves food and isn't picky about what they eat\\
5 & On a secret quest that they can't talk about, except to remind others that they're "on a secret quest"\\
6 & Afraid of something commonplace, such as heights, enclosed spaces, or bare hands and feet\\
7 & Carries around a suitcase full of wigs, forks, ashes, glass orbs, left shoes, or something else strange\\
8 & Full of folksy wisdom that usually starts with, "Papa always said..." or "Mama always said..."\\
\end{DndTable}

\subsection{Fey Outlook}
A Fey creature's outlook can be whatever you want it to be, but rarely do Fey ignore the importance of reciprocity, hospitality, and gifts. These concepts are discussed in the sections that follow.

\subsubsection{Reciprocity}
By and large, Fey feel strongly about quid pro quo and balance. If something is taken, then something of equal value must be given, and what a Fey considers to be of equal value is the big question. A Fey might steal a human's beloved pet and leave in its place a brightly painted wooden effigy, or the Fey might take some gold and leave a bundle of bright, yellow buttercups. These exchanges satisfy the agreement of reciprocity, and a Fey who makes them sleeps soundly at night, content that the exchanges were fair.

\subsubsection{Hospitality}
Hospitality is a pillar of Fey society. Treating a visitor to one's home with with courtesy and generosity is important to most Fey, but the visitor must show their Fey host the same courtesy and not act boorishly or demonstrate blatant impropriety. Rudeness breaks the bond of reciprocity and frees a Fey host from the obligation to be hospitable. But each Fey has different ideas about what constitutes rudeness; even an ill-conceived gift to a Fey host might be regarded as an insult and cause a hubbub, if not a revocation of the Fey host's hospitality. A trusty Fey guide can provide invaluable assistance in navigating such delicate situations.

\subsubsection{Gifts}
Connected to both hospitality and reciprocity is the giving of gifts. Fey are avid gift-givers (partially because they like receiving gifts), and their gifts are usually very thoughtful. Good gifts have sentimental value to the giver. For example, a beloved heirloom makes a fine gift to a Fey creature. A throwaway gift is an insult that flies in the face of reciprocity and hospitality.

Gifts provide a kind of balance that many Fey obsess over. If there is a perceived imbalance, a carefully chosen gift can set things right; this is why refusing a gift from a Fey creature can cause them consternation, as they're trying to right a perceived imbalance by giving the gift. But accepting a gift from a Fey can cause problems for the uninitiated traveler, as it can indebt them to the Fey creature: "I gave you a gift, so now you must give me one in return." It can also lead to the formation of an accidental fey contract, as discussed in the next section.

{\par \vspace { 9pt plus 3pt minus 3pt } \noindent  \DndFontTableTitle{Good Gifts to Give Fey} \nopagebreak }

\begin{description}
\item Family ring
\item Talisman sacred to your druid circle
\item Favorite hat
\item Piece of art that you made
\item Trinket from the Material Plane
\item Favorite dessert recipe
\item Sensational, applause-worthy performance
\item Haircut or bath
\end{description}

{\par \vspace { 9pt plus 3pt minus 3pt } \noindent  \DndFontTableTitle{Bad Gifts to Give Fey} \nopagebreak }

\begin{description}
\item Item that has outlived its usefulness
\item Item you intended to get rid of
\item Half-hearted performance
\item Cursed, evil, or broken item
\end{description}

\subsection{Fey Contracts}
A fey contract is formed when a creature receives a gift (or the promise of a gift) from a Fey and is expected to give the Fey something in return. The gift can be almost anything, and the contract forms as soon as the gift is received.

\subsubsection{Accidental Fey Contracts}
A creature can accidentally stumble into a fey contract in a number of ways. The one bit of good news for the novice Feywild explorer is that most fey contracts can be broken with a \textit{remove curse} spell or similar magic. Only the most powerful fey contracts are hard to break—ones woven by ancient hags, the queens of the Seelie and Unseelie Courts, and other powerful archfey. Such contracts usually require a \textit{wish} spell or an elaborate ritual to negate.

Here are a few ways one might become unwittingly bound to a fey contract:


\begin{itemize}
\item Accepting a gift from a Fey (prompting the Fey to expect something of perceived equal value in return)
\item Stealing something from a Fey (creating a metaphysical imbalance that must be rectified)
\item Taking the life of a creature that made a contract with a Fey (thereby inheriting the creature's debt to that Fey)
\end{itemize}

\subparagraph{Accepting Gifts}
Some Feywild guides recommend never accepting gifts from a Fey and, more importantly, never expressing thanks.

To accept a gift from a Fey is to enter into a contract with it, especially if the gift is received with gratitude. Effusive thanks increase the gift's perceived value, and the Fey will expect something more in return.

\subparagraph{Stealing from a Fey}
Even if a Fey creature is unaware something has been stolen from them, they sense that they have been deprived of something. This nagging sense doesn't go away until the Fey figures out what they lost and who has the stolen item. Moreover, the Fey might not want the stolen thing back, but rather something of equal or greater perceived value.

\subparagraph{Taking a Life in Debt}
Before taking the life of a creature in the Feywild, a wise individual makes sure that creature has no outstanding debts to Fey. Any Fey the creature was indebted to look to the killer to make good on the creature's unfulfilled debts.

Fey are understandably cautious when collecting debts from a person who is prone to violence. They usually make their demands at times when it's safest for them to do so, such as when the killer is bathing or in a crowded place.

\subsubsection{Making a Contract}
Fey contracts can be divided into two categories: greater contracts and lesser contracts. Greater contracts are made with archfey, ancient hags, and other powerful Fey spellcasters. Lesser contracts are made with Fey of all other sorts. Here are some examples of gifts Fey can bestow as part of a greater or lesser contract:

{\par \vspace { 9pt plus 3pt minus 3pt } \noindent  \DndFontTableTitle{Greater Contract Gifts} \nopagebreak }

\begin{description}
\item[Audience.]
You and your companions gain a private audience with the Summer Queen, the Queen of Air and Darkness, or both.
\item[Major.]
You gain a very rare magic item that is yours to keep, or a legendary magic item for 5d12 days.
\item[Safety.]
One creature that previously regarded you as an enemy no longer remembers you at all.
\item[Time.]
You and your companions can return to the Material Plane up to fifty years from now without having aged a day.
\item[Title.]
You gain an important title and all the rewards that come with it (see "Marks of Prestige" in the Dungeon Master's Guide).
\item[Wealth.]
You receive up to 50,000 gp worth of coins, jewelry, or property.
\end{description}

{\par \vspace { 9pt plus 3pt minus 3pt } \noindent  \DndFontTableTitle{Lesser Contract Gifts} \nopagebreak }

\begin{description}
\item[Charm.]
You gain a charm of your choice, subject to the DM's approval (see "Supernatural Gifts" in the Dungeon Master's Guide).
\item[Fey Kinship.]
For 5d6 days, you gain either the Fey Ancestry trait common among elves or the Speak with Small Beasts trait common among forest gnomes.
\item[Guide.]
You and your companions receive help from a guide who can guarantee safe passage through a particular region of the Feywild (such as a Domain of Delight).
\item[Inspiration.]
You gain inspiration each day at dawn for 2d4 days.
\item[Invitation.]
You receive an invitation to the Summer Court or the Gloaming Court (though this invitation does not guarantee an audience with the Summer Queen or the Queen of Air and Darkness).
\item[Minor Magic Item.]
You gain a common magic item that is yours to keep, or an uncommon magic item for 5d6 days.
\item[Proficiency.]
You gain proficiency in a skill of your choice for 2d4 days.
\item[Spellcasting.]
You gain the ability to cast a spell of 4th level or lower once, without material components. You choose the spell, and your spellcasting ability for it is Charisma.
\end{description}

\subsubsection{Price of a Contract}
If you're not sure what a Fey expects to receive for a gift they bestow, roll on the appropriate Fey Desires table.

To collect something intangible, such as a creature's singing voice or the color in its eyes, a Fey must tap into the magic of the Feywild. In other words, it's the magic of the Feywild, not the Fey creature, that allows a character to claim what would otherwise be impossible to obtain. For this reason, a Fey can't claim such a thing unless they and the creature with which they made the contract are both in the Feywild.

\begin{DndTable}[header=Fey Desires for Greater Contracts]{cX}

d8 & What the Fey Wants\\
1 & One of your fingers\\
2 & To take the next child born in your family and raise the child in the Feywild\\
3 & The completion of three quests, each of which you must agree to before the contract is formed\\
4 & Your everlasting fealty\\
5 & The ruin or demise of the Fey's sworn enemy\\
6 & A precious object (such as a rare, very rare, or legendary magic item) that was stolen by or belongs to the Fey's sworn enemy\\
7 & An art object that is deemed priceless, such as a famous painting that hangs in a museum\\
8 & Your youth (which the Fey can harvest and bestow on another creature)\\
\end{DndTable}

\begin{DndTable}[header=Fey Desires for Lesser Contracts]{cX}

d8 & What the Fey Wants\\
1 & Your singing voice\\
2 & A trinket that carries great sentimental value\\
3 & Companionship (you remain in the Fey's company for an agreed-upon period of time)\\
4 & The color in your eyes\\
5 & The memory of your first kiss\\
6 & The spring in your step or the sparkle in your eyes\\
7 & A lock of your hair\\
8 & Your name (requiring you to choose a new one for yourself)\\
\end{DndTable}

\subsubsection{Breaking a Contract}
The Feywild can punish a creature for breaking a fey contract, but the creature must be on the plane to be affected. The penalty imposed on a creature who breaks a lesser contract can be removed by any magic that ends a curse; a \textit{wish} spell is needed to remove the penalty for breaking a greater contract.

The Breaking a Greater Contract and Breaking a Lesser Contract tables can be used to determine the magical penalty of breaking a fey contract.

\begin{DndTable}[header=Breaking a Greater Contract]{cX}

d8 & Penalty\\
1 & You can't speak or cast spells with verbal components. Whenever you try to speak, you bray like a donkey instead.\\
2–3 & You magically transform into an owl. You retain your languages, your ability to speak, and your mental ability scores (Intelligence, Wisdom, and Charisma). You otherwise have the statistics of an owl.\\
4–5 & You can't attune to magic items. If you are currently attuned to one or more magic items, your attunement to those items ends immediately.\\
6–7 & You are petrified.\\
8 & An iron thorn is magically lodged in your chest. Every day, you can feel it move closer to your heart. If this effect is not ended within 3 days, you die and can't be revived for 5d12 days.\\
\end{DndTable}

\begin{DndTable}[header=Breaking a Lesser Contract]{cX}

d8 & Penalty\\
1–2 & Your appearance becomes more toad-like (warty skin, bulbous eyes, large mouth, and webbed digits). This transformation has no game effects.\\
3 & You smell like swamp gas. No amount of bathing can rid you of this stench.\\
4 & Your shadow does not match your movements, which other creatures find unsettling.\\
5–6 & Your appearance becomes more rat-like (beady eyes, whiskers, small ears, pointy nose, little hands and feet, and a rat's tail). This transformation has no game effects.\\
7 & You cast no reflection. Superstitious folk who notice think you're a vampire.\\
8 & You are constantly surrounded by a small cloud of annoying but harmless flies.\\
\end{DndTable}

\subsection{Fey Curses}
Curses are common punishments among archfey and other powerful Fey creatures. An adventurer might be cursed for any number of reasons, a few of which are listed below:


\begin{itemize}
\item Offending a powerful Feywild denizen
\item Entering a forbidden place
\item Appearing in a fey court without an invitation
\end{itemize}

You can determine the curse's effect by rolling on the Fey Curses table.

\begin{DndTable}[header=Fey Curses]{cX}

d8 & Curse\\
1 & Your ears are magically replaced by a pair of soft, fuzzy donkey ears. Moreover, when you try to speak, you instead bray.\\
2 & You gain 1d3 levels of exhaustion. Until the curse ends, these levels of exhaustion can't be removed.\\
3 & Spells can't restore hit points to you.\\
4 & The sound of pixie laughter fills your head while you are awake, drowning out all other sounds.\\
5 & Anything you try to pick up or hold in your hands slips through your fingers.\\
6 & Moonlight burns your flesh. You take 1d10 radiant damage when you start your turn in moonlight.\\
7 & You are magically transformed into an animated wooden doll that looks like you. Your statistics are the same, but you are a Construct with vulnerability to fire damage, and you don't require air, food, or drink. Items worn or carried by you are unaffected.\\
8 & Whenever you tell a lie, you lose the ability to speak for 1d8 hours.\\
\end{DndTable}

A \textit{remove curse} spell or similar magic is usually enough to end a Fey curse on a creature, but some Fey curses are tenacious and resistant to all magic except a \textit{wish} spell. A creature can also remove such a curse on itself by learning and performing a specific task or ritual, determined by rolling on the Ending the Curse table.

\begin{DndTable}[header=Ending the Curse]{cX}

d12 & Ending the Curse\\
1 & You must carve your name into a tree, whereupon the tree inherits the curse's magic and dies.\\
2 & You must drink the blood of a pixie or sprite while basking in moonlight.\\
3 & You must speak the true name of the creature that cursed you three times in a row.\\
4 & While standing in sunlight on an arched bridge over running water, you must beg the Summer Queen for her "fair blessing."\\
5 & After filling your pockets with fool's gold, you must flap your arms and quack like a duck.\\
6 & You must bury an executioner's hood (a blackcapped mushroom found in the Feywild) in the earth and pour goat's milk over it while whistling.\\
7 & You must be bitten by a faerie dragon. (It's possible any faerie dragon might do, or the faerie dragon might have to be of a particular color.)\\
8 & You must leave a bouquet of eight black roses at the place where you were cursed, or you must give the bouquet to the creature that cursed you.\\
9 & You must persuade a centaur to carry you on its back for eight miles. Before the centaur will do this willingly, it might demand a gift or service in exchange, or the completion of a quest.\\
10 & You must obtain leaves or pinecones from three different species of treants and burn them in a campfire under a full moon while singing a particular campfire song.\\
11 & You must bake a small cake and leave it on the doorstep of a forest gnome's abode. Only when the gnome eats the whole cake does your curse end.\\
12 & You must persuade another creature to willingly take your name. If it does so, it inherits your curse, whereupon you are nameless and must choose a new name for yourself.\\
\end{DndTable}

\subsection{Fey Abodes}
Fey creatures live in abodes they fashion for themselves or repurpose for their needs. Characters might stumble upon these abodes in their Feywild wanderings. Use the Fey Abodes table to randomly determine a Fey creature's lair, or choose an option that works well for the creature in question.

\begin{DndTable}[header=Fey Abodes]{cX}

d20 & Abode\\
1 & A pagoda overgrown with flowering vines that beckon visitors with pleasing scents\\
2 & A crooked stone tower that has a moon-like orb of light circling its rooftop, which causes the tower's shadow to move like a clock hand\\
3 & A decrepit mansion that is partially sunk in the middle of a bog\\
4 & A rocky hill shaped like a sleeping satyr, with its open mouth forming a cave entrance\\
5 & A windmill that walks around on giant crow's feet\\
6 & A crumbling keep on a small island in the middle of a mist-shrouded loch\\
7 & An old farm overgrown with giant pumpkins\\
8 & A treehouse built in the boughs of a treant\\
9 & A tower that used to be the trunk of a giant petrified tree, with rope bridges connecting it to the giant living trees that surround it\\
10 & One or more houses in hollowed-out mushrooms\\
11 & A gingerbread cottage with a frosting-covered roof, frosting icicles, chocolate doors, and gumdrop gardens\\
12 & A musty, web-shrouded stone cottage surrounded by an orchard of awakened apple trees\\
13 & A walking stone colossus with a tower for a head\\
14 & A cave-riddled hill that walks around on giant stone feet\\
15 & A walled garden filled with friendly critters, talking flowers, and grasping vines\\
16 & A well-preserved elven tomb overgrown with moss, decorated with statues, and festooned with bird nests\\
17 & A giant beaver's lodge\\
18 & An inn or hostel carved into the foot of a hill\\
19 & A dragon skull lying in the sand\\
20 & One or more giant rusty helmets that serve as houses, surrounded by an ancient battlefield\\
\end{DndTable}

\chapter{Creating an Archfey}
\DndDropCapLine{A}{}rchfey are among the most powerful beings in the Feywild. Consciously or unconsciously, they transform their Feywild homes into reflections of their desires and complex personalities. Powerful archfey such as the Summer Queen and the Queen of Air and Darkness rule vast domains, and their influence is so great that their courts stretch beyond the borders of their Domains of Delight to the far ends of the Feywild. Other, lesser known archfey rule domains that are minuscule by comparison, but no less wondrous.


\section{Where to Start?}
Most archfey are ancient creatures, and no two are alike. When creating an archfey, here are some good questions to ask yourself:


\begin{itemize}
\item Under what circumstances might adventurers encounter the archfey?
\item Does the archfey have any allegiance to the Summer Court or the Gloaming Court?
\item What traits does the archfey have? Specifically, what do they look like, what is their personality, and what magical powers do they possess?
\end{itemize}

These questions are tackled in the sections that follow. The answers will help you flesh out the archfey's Domain of Delight.

\subsection{Encountering the Archfey}
Think about how and when you expect the adventurers to encounter your archfey. The archfey might infiltrate the party using illusion magic or spy on the characters from afar to discern their intentions before confronting them. The archfey might test the characters' mettle with a series of challenging encounters or puzzles before deeming them fit to bask in the archfey's presence. Conversely, the archfey might be blissfully unaware of the party's intrusion or show no interest in the characters. A shy or reclusive archfey might not want to be disturbed at all.

How and when the archfey crosses paths with the adventurers can be determined by answering the following questions:


\begin{itemize}
\item Does the archfey engage with visitors or shun them?
\item Does the archfey wander their domain, or do they prefer to remain in their lair?
\item Does the archfey visit other domains? Are they involved in the politics of the Fey courts?
\end{itemize}

\subsection{Fey Court Allegiances}
The Summer Queen and the Queen of Air of Darkness are so powerful that other archfey might feel it's in their best interest to ally with one or both of them. By swearing allegiance to either queen, an archfey gains the privilege of audiences with the queen and faint assurances that the queen's court will not meddle in the affairs of the archfey or their domain.

Archfey that do not pledge allegiance to one court or the other can find their domains overrun with Seelie or Unseelie spies or overlapped by the domains of archfey in league with one or both queens.

\subsection{Presentation}
Archfey have one thing in common: they're all Fey. Beyond that, their features vary widely. To randomly determine what your archfey looks like, roll on the Archfey's Presentation table.

\begin{DndTable}[header=Archfey's Presentation]{cX}

d8 & Presentation\\
1 & \textbf{Adorable.} The archfey has taken on the form of something adorable, such as an otter, a fawn, a sugar glider, a wombat, or a beautiful butterfly.\\
2 & \textbf{Ancient.} This archfey appears to be eons old. They might have gnarled skin like teakwood and long, gray hair that looks like elaborate lacework.\\
3 & \textbf{Bizarre.} This archfey looks like something from the realm of dreams. They could take on the appearance of a scintillating ball of light, a floating mask, or a whispering shadow.\\
4 & \textbf{Floral.} Flowers bloom from the archfey's fingertips, and their skin is covered in petals. Their body might have thorns, or they could have vine-like hair covered in sweet-smelling blossoms.\\
5 & \textbf{Fluid.} The archfey can change body type or transform from one creature into another.\\
6 & \textbf{Iridescent.} The archfey floats off the ground and gazes about with glowing eyes, their body surrounded by a nimbus of light and their skin sparkling like sunlight on water.\\
7 & \textbf{Metallic.} The archfey appears to be made of metal. Their face is an exquisitely crafted mask of mithral, silver, or gold, and their body is made of finely tooled metal etched with organic designs or cryptic symbols.\\
8 & \textbf{Monstrous.} This archfey has a monstrous form. For example, they might resemble an enormous spider, a hydra whose heads look like oversized elf heads with fangs, or any other horrific form you can imagine.\\
\end{DndTable}

\subsection{Personality}
If you haven't already settled on a personality for your archfey, you can determine it randomly by rolling on the Archfey's Personality table.

\begin{DndTable}[header=Archfey's Personality]{cX}

d8 & Personality\\
1 & \textbf{Benevolent.} The archfey welcomes all into their domain and does everything in their power to assist those who come to them in true friendship or grave need.\\
2 & \textbf{Covetous.} The archfey hoards things they consider valuable and uses spies to find more of what they covet. The archfey never leaves their lair for fear that their hoard might be plundered.\\
3 & \textbf{Imperious.} The archfey surrounds themself with a court of sycophants and expects visitors to grovel before them. Their mercy is rare, and their superiority complex is legendary.\\
4 & \textbf{Madcap.} The archfey loves to sing, dance, drink, wear silly masks, do cartwheels, and throw parties where everyone is free to cavort and revel how they wish.\\
5 & \textbf{Mercurial.} The archfey has two personalities (choose two other personalities from this table). They shift back and forth between these personalities at certain times or when certain conditions are met.\\
6 & \textbf{Mischievous.} The archfey is a trickster who likes to confound and frustrate visitors to their domain and who surrounds themself with sly creatures such as boggles, pixies, foxes, magpies, and crows.\\
7 & \textbf{Reclusive.} The archfey prefers to be left alone and seldom, if ever, leaves their domain. They might adopt disguises or turn invisible when they travel, and they might use magic to conceal their lair.\\
8 & \textbf{Wild.} The archfey has the demeanor of a wild animal and shuns the trappings of society.\\
\end{DndTable}

\subsection{Obsessions}
Obsessions are common among archfey. An archfey might be obsessed with a painting of themself and spend days, months, or years staring at it. An archfey might obsess over a perceived injustice and work tirelessly to correct it. An archfey might obsess over their collection of acorns, believing the acorns are keys that unlock the secrets of the multiverse.

If you want your archfey to have an obsession, you can determine it randomly by rolling on the Archfey Obsession table.

\begin{DndTable}[header=Archfey Obsession]{cX}

d8 & Obsession\\
1 & \textbf{Beauty.} The archfey does everything in their power to rid their domain of that which they consider ugly. What they consider beautiful is not always clear, however.\\
2 & \textbf{Color.} This archfey is obsessed with a particular color and puts their subjects to work making sure the archfey's favorite color is predominant. Wearing another color in the domain could be seen as laughable, obnoxious, or downright rude.\\
3 & \textbf{Etiquette.} The rules in the archfey's domain change as they learn about some new custom that is popular in the Summer Court. The archfey requires all creatures in their domain to respect the new rules of etiquette.\\
4 & \textbf{Magic.} Magic in all its forms fascinates the archfey. They collect magic items and spellbooks, and they enjoy seeing demonstrations of magic from visitors.\\
5 & Material \textbf{Plane.} The archfey is obsessed with creatures and objects from the Material Plane. There is a scent, an energy, and a uniqueness to them that fascinates the archfey.\\
6 & \textbf{Monsters.} The archfey has a menagerie of monsters and is obsessed with finding marvelous new creatures to add to its collection.\\
7 & Rule of \textbf{Three.} The archfey is obsessed with the Rule of Three. They look for patterns, deeper meanings, and ill omens in things that occur in threes. The archfey's obsession is reflected throughout their domain, where things seem to naturally come in threes.\\
8 & \textbf{Stories.} The archfey wants nothing more than to listen to stories told by creatures from faraway places. They love to gather lore so they can build their own internal world of imagination or transform their domain with ideas from other worlds. Their lair is festooned with books, which give the archfey inspiration.\\
\end{DndTable}

\subsection{Signature Magic}
Almost every archfey has a signature magical power, and this power might be tied to their appearance or personality; for example, a metallic archfey might have the power to corrode or heat metal, while a madcap archfey might have the power to cause wild magic effects similar to those described the Wild Magic Surge table in the "Sorcerer" section of the Player's Handbook. The Signature Magic table presents some other possibilities.

\begin{DndTable}[header=Signature Magic]{cX}

d8 & Magic\\
1 & \textbf{Arcane Antlers.} The archfey has antlers made of metal, crystal, bone, or some other substance, which allow the archfey to raise the dead. They shed their antlers and grow new ones every hundred years, prompting thieves to sneak into the domain to acquire the shed antlers.\\
2 & \textbf{Compelling Voice.} The archfey has a magical voice that can charm creatures of a certain type (such as Beasts or Fey). This voice might be as soft as a whisper or loud like thunder.\\
3 & \textbf{Dreamwalk.} When they sleep, the archfey projects a ghost-like version of themself that can leave their body and go anywhere in their domain.\\
4 & \textbf{Hand of Knowing.} One of the archfey's hands blurs as though it were in multiple places at once. The archfey can learn the entire history of a creature or object by touching it with that hand.\\
5 & \textbf{Health Aura.} The archfey's presence rids creatures and vegetation of disease. Healthy vegetation yields an overabundance of flowers and fruit.\\
6 & \textbf{Scrying Eye.} This archfey has an ornately crafted, removable eye that has the power of truesight (as described in the Monster Manual). The archfey can see through this orb at all times. If the detached eye is destroyed, it rematerializes undamaged in the archfey's empty eye socket.\\
7 & \textbf{Simulacra.} The archfey can make magical copies of themself and other creatures. These simulacra are similar to those created by the \textit{simulacrum} spell.\\
8 & \textbf{Time Distortion.} For each minute spent in the archfey's presence, an hour passes elsewhere.\\
\end{DndTable}

\subsection{Magical Gifts}
Some archfey have the power to give magical gifts to those they deem worthy. A gift could be anything from a trinket that grants safe passage through the archfey's domain to a magic item that can be used to defeat a powerful monster. As the DM, you decide whether such a gift is warranted. To randomly determine what kind of gift your archfey might bestow on a character, roll on the Magical Gifts table.

\begin{DndTable}[header=Magical Gifts]{cX}

d8 & Gift\\
1 & \textbf{Charm.} The archfey bestows a charm of the DM's choice (see "Supernatural Gifts" in the Dungeon Master's Guide).\\
2 & \textbf{Guide.} The archfey summons a magical guide that takes the form of a faerie dragon, a sprite, or some other creature. The guide is a harmless figment that can't be damaged, and it knows its way around the Feywild. It can't leave the Feywild and disappears after eight days.\\
3 & \textbf{Lore.} With a touch, the archfey magically imparts useful knowledge about a particular subject of the archfey's choice. This knowledge might be permanent, or it might fade over time.\\
4 & Magic \textbf{Item.} The archfey bestows a useful magic item (such as a crystal ball or an oathbow) but warns that the item will vanish after a specified period of time, which it does.\\
5 & \textbf{Refuge.} The archfey creates a wooded glen, a dome of brambles, a silken cocoon, or some other kind of refuge. Creatures hostile toward the gift's recipient can't enter this refuge without the recipient's consent. After a specified period of time, the refuge vanishes.\\
6 & \textbf{Resurrection.} The archfey offers one free casting of the \textit{true resurrection} spell, which they can cast as an action without material components once per century.\\
7 & \textbf{Training.} With a touch, the archfey grants the benefit of months of special training (see "Marks of Prestige" in the Dungeon Master's Guide for different benefits of training).\\
8 & \textbf{Transformation.} The archfey offers one free casting of the \textit{true polymorph} spell, which they can cast without material components. The spell's effect can be made permanent, if you wish.\\
\end{DndTable}

\section{Archfey Statistics}
If your campaign includes an archfey whom the characters might end up fighting, create a stat block for the archfey or modify some other creature's stat block to suit your needs. The Dungeon Master's Guide contains guidance for modifying existing stat blocks.

Alternatively, choose a stat block of CR 5 or higher that has one of the following creature types: Beast, Elemental, Fey, Giant, Humanoid, or Plant. Change the creature's type to Fey, make its Intelligence at least 10 (+0), and give it two languages if doesn't speak any languages. Also give the creature one or both of the following traits:

\subparagraph{Fey Rebirth}
If the archfey dies in its Domain of Delight, it revives with all its hit points 1d4 days later in a safe location in that domain.

\subparagraph{Legendary Resistance (3/Day)}
If the archfey fails a saving throw, it can choose to succeed instead.

\chapter{Creating a Domain of Delight}
\DndDropCapLine{O}{}nce your archfey is fleshed out, you can create the Domain of Delight that serves as the archfey's home in the Feywild. Generally, the size of a Domain of Delight is a good determiner of an archfey's power, but an archfey's domain can be as big or as small as you want it to be.


\section{Geographical Features}
A small Domain of Delight might have only one prominent geographical feature, which you can determine randomly by rolling on the Geographical Features table. A large Domain of Delight can have multiple geographical features, which you can determine by rolling more than once on the table.

\begin{DndTable}[header=Geographical Features]{cX}

d8 & Feature\\
1 & A colossal, gnarled tree, its trunk riddled with passageways, halls, chambers, and staircases\\
2 & A valley filled with grassy meadows and groves of trees along a running river, nestled between two snow-capped mountains that look like the points of a crescent moon\\
3 & A crystal castle that sings in sunlight or moonlight\\
4 & A vast thicket of thick roots, thorny vines, and sinuous creepers that weave together to form long tunnels, grand hallways, and enormous domes\\
5 & A complex of caves connected by passageways formed by Fey magic to create grand galleries and soaring vaults\\
6 & Rolling farmland that yields healthy, bountiful produce of amazing size\\
7 & A forest of tower-sized mushrooms, or a forest of tiny mushrooms that characters must shrink down to explore\\
8 & An island that is actually a colossal sleeping turtle\\
\end{DndTable}

\section{Domain Theme}
A Domain of Delight might have a prevailing theme, which you can determine randomly by rolling on the Domain Theme table.

\begin{DndTable}[header=Domain Theme]{cX}

d8 & Theme\\
1 & \textbf{Adversarial.} Adversarial personalities are common here, even among the local wildlife. Animated trees swat at passersby, sentient flowers scoff and sneer, and local Fey are grumpy and uncooperative.\\
2 & \textbf{Crystalline.} This domain contains a profusion of crystal formations and outgrowths, as well as structures and trees made of crystal. Some of the crystals might generate light, music, or both.\\
3 & \textbf{Gloomy.} Everything here has a gloomy cast. Tree branches look like skeletal fingers, night creatures slink and flutter about, and skull-like faces appear on plants, stones, and sprites' wings.\\
4 & \textbf{Inquisitive.} Creatures in this domain are nosy and hungry for gossip. Animated trees use their branches to pick through backpacks when travelers aren't looking; pixies spy on strangers; and birds eavesdrop on conversations, repeating what they've heard to their masters.\\
5 & \textbf{Kaleidoscopic.} This domain is a riot of everchanging colors. Plants come in a dazzling array of colors, flowers change color to suit the prevailing mood, and Fey wear costumes of scintillating hues.\\
6 & \textbf{Lackadaisical.} The creatures here are easygoing. No one is in a hurry, and no task is greeted with a sense of urgency.\\
7 & \textbf{Musical.} This domain abounds with music. Frogs, toads, and insects break into multipart harmonies that sometimes crescendo into a symphonic rapture of melodious tweets, trills, chirps, and whistles. Fey in this domain are just as musically inclined.\\
8 & \textbf{Poisonous.} Almost every plant here has spines or thorns that cause anything from itchy welts to damage or out-of-body experiences. Other sources of poison might be present as well (for example, toxic bogs or giant, poisonous Beasts).\\
\end{DndTable}

\section{Coming and Going}
Travelers who don't have access to \textit{plane shift} spells or similar magic must rely on other means to travel to and from a Domain of Delight. Fey crossings are the most common means of transit between the Material Plane and the Feywild; more information about them can be found in the Dungeon Master's Guide.

An archfey is usually aware of any Fey crossings that allow creatures to enter and leave the archfey's domain and may protect these Fey crossings or leave them unguarded. Some Fey crossings operate only at certain times or when certain conditions are met. The Fey Crossings table allows you to randomly determine the characteristics of a fey crossing that allows passage to and from a Domain of Delight. Your Domain of Delight can have as few or as many Fey crossings as you like, and they need not all be the same in appearance or function.

Geographical features on one side of a fey crossing tend to be similar to those on the other side, to the extent that some travelers might be unaware they've left one plane of existence and entered another. For example, if one stumbles upon a fey crossing in a cave, a similar cave might lie on the other side of that fey crossing, as if the caves were near-perfect reflections of each other.

\begin{DndTable}[header=Fey Crossings]{cX}

d8 & Fey Crossing\\
1 & \textbf{Altar.} This stone altar might be found atop a lonely plateau, behind a waterfall, inside a cave, or in some other remote location. When certain stars align and the proper ritual is performed, the fey crossing appears as a portal above the altar.\\
2 & \textbf{Crystal Cave.} This cave is filled with natural crystalline formations and resembles a glittering cathedral. When moonlight shines into the cave at night, it causes a shimmering pool or spring in the center of the cave to transform into a fey crossing.\\
3 & \textbf{Hollow Tree.} The fey crossing lies within the hollow interior of an enormous tree. The crossing is open to anyone who carries a leaf or acorn from the tree.\\
4 & \textbf{Maze.} The fey crossing lies at the center of a maze that might change configuration from time to time. To use the crossing, one must walk or run from the maze's entrance to its center. Creatures who circumvent the maze using magic or flight find themselves unable to use the crossing.\\
5 & \textbf{Mushroom Circle.} This fey crossing appears as a circle of mushrooms. The size of the mushrooms and the width of the circle can vary, but activating the fey crossing requires that a particular act be performed inside it or a particular object or creature be present.\\
6 & \textbf{Ruined Tower.} The ruins of an ancient elven tower house a fey crossing. For three nights each year, a ghostly image of the intact tower appears above the ruined foundation, during which time the fey crossing appears as a shimmering doorway that creatures can pass through.\\
7 & \textbf{Stone Bridge.} A Fey creature hides under this ancient stone bridge. In exchange for treasure or some other gift, this Fey can use its innate magic to create a portal atop the bridge that serves as a fey crossing.\\
8 & \textbf{Stone Circle.} This circle of stone megaliths contains a fey crossing that is active during equinoxes and solstices.\\
\end{DndTable}

\section{Travel in the Domain}
Visitors to a Domain of Delight are often confounded by the Feywild's power to distort distance, time, and reality. A traveler might see a hilltop that appears to be a mile away and march toward it, only to find the hilltop getting farther away. Similarly, a wanderer who asks for directions might be told to follow a particular trail, only to discover the trail changes course depending on the mood of those who travel along it. Determining how far one has traveled in a particular amount of time can be a headscratcher.

\subsection{Magical Shortcuts}
Domain denizens and guides know some of the tricks to navigating a Domain of Delight, including magical shortcuts. A shortcut can also be found by accident or learned in some other way. For example, a traveler might wake from a vivid dream or a trance with knowledge of a shortcut that fades from memory after an hour.

To take advantage of a magical shortcut, one must spend 1 minute performing a routine; the Shortcut Routines table provides examples. At the end of this routine, the performer and any creatures traveling with them find themselves closer to their intended destination. How much closer is up to the DM, but a shortcut routine usually reduces the travel time to the intended destination by at least half.

A shortcut routine allows for safe, expeditious travel from one location in the Feywild to a specific destination in the Feywild, and that's all. For example, it might expedite travel from the edge of an archfey's domain to the entrance of the archfey's lair but not the other way around, and it wouldn't shorten the travel time to other locations in the domain.

\begin{DndTable}[header=Shortcut Routines]{cX}

d8 & Routine\\
1 & Picking petals from a flower while walking toward a beckoning sunset\\
2 & Playing "Ode to the Summer Queen" on a \textit{lute} while skipping counterclockwise around a tree\\
3 & Eating a mushroom while gazing at one's own reflection in a clear pool or stream of water\\
4 & Doing a headstand in a ring of stones\\
5 & Staring at firelight while playing a \textit{flute} or singing about a fond memory\\
6 & Sitting by a creek while lost in childlike wonder\\
7 & Whistling a tune while walking backward\\
8 & Holding hands and dancing around a tree at twilight\\
\end{DndTable}

\section{Drama}
You can bring your Domain of Delight to life by adding drama. When the characters show up, they might become swept up in this drama, perhaps even embroiled in intrigue. The drama can be anything from a petty squabble between Fey settlements to a full-blown fomorian incursion. To determine a worthy conflict, you can roll on the Domain Drama table or choose an entry that you like.

\begin{DndTable}[header=Domain Drama]{cX}

d8 & Drama\\
1 & A cherished object or favorite pet belonging to the domain's archfey ruler has gone missing.\\
2 & Fomorians (or other hostile creatures) have emerged from a cave and are causing a ruckus.\\
3 & The domain's archfey ruler recently turned down an invitation to the Gloaming Court, offending the Queen of Air and Darkness, whose spies are now sowing discord throughout the domain.\\
4 & The domain's archfey has fallen into a magical slumber. Others are looking for someone they can trust to find a way to awaken the archfey.\\
5 & The vegetation is sick in part of the domain, and the contagion is spreading. The domain's archfey needs help finding the cause of the sickness.\\
6 & Part of the domain overlaps the domain of another archfey, and the Fey living in that region are tired of answering to two masters. They won't be happy until one archfey yields to the other.\\
7 & A sleeping knight in silver plate mail bearing no distinguishing symbols has been discovered in a hollow tree trunk. Local Fey wonder who she is and if they should revive her or not.\\
8 & An arranged wedding was supposed to unite two feuding villages, but one of the grooms has disappeared. Everyone is searching for him in the hope that the wedding will take place, bringing peace to the domain.\\
\end{DndTable}

\section{Weird Magic}
Magical forces, strange interlopers, and tragic events can alter a smaller region within a Domain of Delight. While the flora, fauna, structures, and inhabitants might remain unaffected, the region's innate character takes on new qualities.

When you wish to include such a region, roll on the Weird Magic Effects table to determine its effects, or choose an entry you like. These effects occur whenever you please, at the time each description suggests, or under one or more of the following circumstances:


\begin{itemize}
\item Soon after the party first enters the region
\item When a creature loses more than half its hit points
\item When a creature casts a spell of 1st level or higher
\item When a creature activates a magic item
\item When a creature makes an exceptionally loud noise or otherwise attracts attention
\item When the party spends at least 30 minutes in the same region
\end{itemize}

\begin{DndTable}[header=Weird Magic Effects]{cX}

d8 & Effect\\
1 & Time visibly moves more quickly outside the region than inside it. For every minute the party spends inside the region, one hour passes outside the region.\\
2 & Flowers not only are abundant in the region but also have faces and limited sentience, allowing them to communicate with creatures in the Elvish and Sylvan tongues. They know what has transpired in the region within the past day and happilyshare this information.\\
3 & Each character in the region gains the benefit of a \textit{barkskin} spell that lasts for 8 hours.\\
4 & Each character who spends 1 hour in the region undergoes a magical transformation, gaining fur, a tail, and large ears reminiscent of a donkey's, goat's, or llama's. This effect ends on a character 3d8 hours after the character leaves the region. Any magic that ends a curse also restores the character's normal appearance.\\
5 & Harmless fireflies gather and flit around one character in the region, sharing that character's space for 1 hour. The fireflies generate bright light out to a range of 10 feet and dim light for an additional 10 feet. While surrounded by the fireflies, the character feels a profound sense of purpose and gains the benefit of the \textit{bless} spell.\\
6 & Creatures in the region can't leave it and find themselves covering the same ground over and over. By the time they realize this, 3d8 hours have passed, during which they have made no progress in their effort to leave. The effect then ends.\\
7 & One character in the region sees a wide, grinning mouth that floats in midair, 10 feet away from them. No other creature can see the grinning mouth. In a language the character understands, the mouth says, "I'll answer three questions you put to me. Noes and yesses, I promise thee." The character can ask it three questions as if using the \textit{commune} spell.\\
8 & One character in the region triggers a wild magic effect that is determined by rolling on the Wild Magic Surge table in the "Sorcerer" section of the Player's Handbook. If the effect duplicates a spell, the character's spellcasting ability for that spell is Charisma.\\
\end{DndTable}

\begin{DndSidebar}[float=!b]{Awakened Beasts}
Awakened Beasts (Beasts that have received the benefits of an \textit{awaken} spell) are common in the Feywild, where they often serve as spies and companions.



\end{DndSidebar}

\section{Domain Denizens}
The Feywild Creatures table can help you populate a Domain of Delight. Stat blocks for the listed creatures appear in the \textit{Monster Manual} or The Wild Beyond the Witchlight.

Just because a creature doesn't appear on the Feywild Creatures table doesn't mean that it can't appear in a Domain of Delight. Any D\&D creature can find its way to the Feywild, and sometimes it's fun to surprise players with creatures their characters would not expect to meet in the Feywild, such as a band of githyanki or a vampire.

\begin{DndTable}[header=Feywild Creatures]{cX}

CR & Creatures\\
0 & \textbf{Awakened shrub}, \textbf{baboon}, \textbf{badger}, \textbf{campestri}, \textbf{cat}, \textbf{crab}, \textbf{deer}, \textbf{eagle}, \textbf{frog}, \textbf{giant fire beetle}, \textbf{goat}, \textbf{hawk}, \textbf{hyena}, \textbf{jackal}, \textbf{lizard}, \textbf{myconid sprout}, \textbf{octopus}, \textbf{owl}, \textbf{rat}, \textbf{raven}, \textbf{scorpion}, \textbf{sea horse}, \textbf{shrieker}, \textbf{spider}, \textbf{vulture}, \textbf{weasel}\\
1/8 & \textbf{Bandit}, \textbf{blood hawk}, \textbf{boggle}, \textbf{brigganock}, \textbf{camel}, \textbf{displacer beast kitten}, \textbf{flying snake}, \textbf{giant crab}, \textbf{giant rat}, \textbf{giant weasel}, \textbf{guard}, \textbf{harengon brigand}, \textbf{mastiff}, \textbf{merfolk}, \textbf{mule}, \textbf{poisonous snake}, \textbf{pony}, \textbf{stirge}, \textbf{twig blight}\\
1/4 & \textbf{Blink dog}, \textbf{boar}, \textbf{bullywug}, \textbf{constrictor snake}, \textbf{draft horse}, \textbf{drow}, \textbf{elk}, \textbf{flying sword}, \textbf{giant badger}, giant crane (\textbf{pteranodon}), \textbf{giant frog}, \textbf{giant lizard}, \textbf{giant owl}, \textbf{giant poisonous snake}, \textbf{giant snail}, \textbf{giant wolf spider}, \textbf{goblin}, \textbf{harengon sniper}, \textbf{needle blight}, \textbf{panther}, \textbf{pixie}, \textbf{pseudodragon}, \textbf{riding horse}, \textbf{sprite}, \textbf{swarm of ravens}, \textbf{violet fungus}, \textbf{wolf}\\
1/2 & \textbf{Ape}, \textbf{black bear}, \textbf{cockatrice}, \textbf{crocodile}, \textbf{darkling}, \textbf{giant dragonfly}, \textbf{giant goat}, \textbf{giant sea horse}, \textbf{giant wasp}, \textbf{hobgoblin}, \textbf{myconid adult}, \textbf{satyr}, \textbf{scout}, \textbf{swarm of insects}, \textbf{vine blight}, \textbf{warhorse}, \textbf{worg}\\
1 & \textbf{Animated armor}, \textbf{brown bear}, \textbf{bugbear}, \textbf{dryad}, faerie dragon (young), \textbf{giant eagle}, \textbf{giant hyena}, \textbf{giant octopus}, \textbf{giant spider}, \textbf{giant toad}, \textbf{giant vulture}, \textbf{goblin boss}, half-ogre, \textbf{harpy}, \textbf{hippogriff}, \textbf{lion}, \textbf{quickling}, \textbf{scarecrow}, \textbf{spy}, \textbf{swarm of campestris}, \textbf{tiger}\\
2 & \textbf{Awakened tree}, \textbf{bandit captain}, \textbf{berserker}, \textbf{centaur}, \textbf{darkling elder}, \textbf{druid}, \textbf{ettercap}, faerie dragon (old), \textbf{giant boar}, \textbf{giant constrictor snake}, \textbf{giant elk}, \textbf{griffon}, \textbf{merrow}, \textbf{myconid sovereign}, \textbf{ogre}, \textbf{peryton}, \textbf{sea hag}, \textbf{swarm of poisonous snakes}, \textbf{treant sapling}, \textbf{will-o'-wisp}\\
3 & \textbf{Basilisk}, \textbf{bugbear chief}, \textbf{bullywug knight}, \textbf{bullywug royal}, \textbf{displacer beast}, \textbf{giant scorpion}, \textbf{green hag}, \textbf{hobgoblin captain}, \textbf{knight}, \textbf{manticore}, \textbf{minotaur}, \textbf{owlbear}, \textbf{phase spider}, \textbf{redcap}, \textbf{veteran}, \textbf{water weird}, \textbf{winter wolf}, \textbf{yeti}\\
4 & \textbf{Banshee}, \textbf{elephant}, \textbf{ettin}, \textbf{sea hag} (in coven)\\
5 & \textbf{Bulette}, \textbf{drow elite warrior}, \textbf{giant crocodile}, \textbf{giant shark}, \textbf{gorgon}, \textbf{green hag} (in coven), \textbf{hill giant}, \textbf{shambling mound}, \textbf{troll}, \textbf{unicorn}\\
6 & \textbf{Chimera}, \textbf{cyclops}, \textbf{galeb duhr}, \textbf{hobgoblin warlord}, \textbf{mage}, \textbf{mammoth}, \textbf{medusa}, \textbf{wyvern}\\
7 & \textbf{Drow mage}, \textbf{giant ape}, \textbf{korred}, \textbf{oni}, \textbf{stone giant}, \textbf{tree blight}\\
8 & \textbf{Fomorian}, \textbf{frost giant}, \textbf{hydra}\\
9 & \textbf{Abominable yeti}, \textbf{cloud giant}, \textbf{fire giant}, \textbf{treant}\\
11 & \textbf{Behir}, \textbf{roc}\\
12 & \textbf{Archmage}\\
13 & \textbf{Jabberwock}, \textbf{storm giant}\\
17 & \textbf{Dragon turtle}\\
\end{DndTable}

\subsection{Domain Visitors}
As they explore a domain, the characters might encounter other visitors, including friends and foes of the domain's archfey, planar travelers, and lost souls trying to find a way home. Use the Domain Visitors table to generate encounters with such folk.

\begin{DndTable}[header=Domain Visitors]{cX}

d8 & Visitor\\
1 & A ranger (use the \textbf{scout} stat block) searching for a missing companion, lover, or long-lost ancestor\\
2 & A \textbf{druid} searching for herbs, mushrooms, or seeds found only in the Feywild\\
3 & A \textbf{mage} or an \textbf{archmage} who wants to serve the Summer Queen or the Queen of Air and Darkness but doesn't know how or where to find her\\
4 & An elf (use the \textbf{noble} stat block) from the Material Plane who is researching their Fey ancestry\\
5 & A hunter (use the \textbf{gladiator} or \textbf{scout} stat block) who was lured into the Feywild by their prey\\
6 & A \textbf{githzerai zerth} studying the chaotic energy of the Feywild and the magic of Fey creatures\\
7 & A would-be warlock (use the \textbf{cultist} or \textbf{cult fanatic} stat block) seeking an audience with the domain's archfey in the hope of gaining its patronage\\
8 & A friendly \textbf{arcanaloth} using \textit{alter self} spells to assume Humanoid form as it searches the Feywild for one of the Books of Keeping\\
\end{DndTable}

\chapter{Putting It All Together}
\textbf{Yarnspinner} is an archfey created using the tables in part 2. His domain, Fablerise, was fleshed out using the tables in part 3. Fablerise is a small Domain of Delight—roughly three square miles nestled in a forest called Thither, which is part of a much larger Domain of Delight called Prismeer. \textbf{Yarnspinner} only recently arrived in Thither and is unaware that he is intruding upon another archfey's territory (see The Wild Beyond the Witchlight for more information about Prismeer and its archfey ruler).

\chapter{Adventure Aids}
\DndDropCapLine{T}{}he following pages provide easy-to-use versions of appendices D and E from The Wild Beyond the Witchlight, as well as the adventure's Story Tracker.



\begin{itemize}
\item See the Roleplaying Cards deck entry for individual cards, or Appendix D of Wild Beyond the Witchlight for printable assets.
\item See the Stagefright's Lines table to roll, or Appendix E of Wild Beyond the Witchlight for printable assets.
\item See the Story Tracker Appendix of Wild Beyond the Witchlight for printable story tracker assets.
\end{itemize}

\chapter{Credits}

\begin{description}
\begin{multicols}{2}

\begin{description}
\item[Writer.]
\DndDropCapLine{A}{}dam Lee


\item[Art Director.]
Kate Irwin

\item[Developers.]
Jeremy Crawford, Ben Petrisor, Christopher Perkins

\item[Editors.]
Judy Bauer, Jeremy Crawford, Christopher Perkins

\item[Proofreader.]
Ari Levitch

\item[Graphic Designer.]
Trish Yochum

\item[Consultants.]
Shiaw-Ling Lai, Omar Ramadan-Santiago

\item[Interior Illustrators.]
Helder Almeida, Mark Behm, Zoltan Boros, Alayna Danner, Toma Feizo Gas, Lars Grant-West, Sam Keiser, Livia Prima, Ned Rogers, Magali Villeneuve

\item[Concept Art Director.]
Shawn Wood

\item[Concept Illustrators.]
Jedd Chevier, Daarken, Toma Feizo Gas, Titus Lunter, April Prime, Ilya Shkipin, Cory Trego-Erdner, Shawn Wood, Kieran Yanner

\item[Imaging Technician.]
Kevin Yee

\end{description}

\section{D\&D Studio}

\begin{description}
\item[Executive Producer.]
Ray Winninger

\item[Game Design Architects.]
Jeremy Crawford, Christopher Perkins

\item[Design Manager.]
Steve Scott

\item[Design Department.]
Sydney Adams, Judy Bauer, Makenzie De Armas, Dan Dillon, Amanda Hamon, Ari Levitch, Ben Petrisor, Taymoor Rehman, F. Wesley Schneider, James Wyatt

\item[Art Department Manager.]
Richard Whitters

\item[Principal Art Director.]
Kate Irwin

\item[Art Department.]
Trystan Falcone, Emi Tanji, Shawn Wood, Trish Yochum

\item[Senior Producer.]
Dan Tovar

\item[Producers.]
Bill Benham, Robert Hawkey, Lea Heleotis

\item[Director of Product Management.]
Liz Schuh

\item[Product Managers.]
Natalie Egan, Chris Lindsay, Hilary Ross, Chris Tulach

\end{description}

\section{Marketing}

\begin{description}
\item[Director of Global Brand Marketing.]
Brian Perry

\item[Global Brand Manager.]
Shelly Mazzanoble

\item[Senior Marketing Communications Manager.]
Greg Tito

\item[Community Manager.]
Brandy Camel

\end{description}

\section{D\&D Beyond}

\begin{description}
\item[D\&D Beyond Product Manager.]
Patrick Backmann

\item[D\&D Beyond Digital Design Team.]
Jay Jani, Adam Walton, Cameron Powell, Joseph Keen

\end{description}

\end{multicols}

\end{description}

\end{document}
